\section{Lập trình Di truyền}
\label{sec:gp}

Lập trình Di truyền do John Koza phát triển
từ cuối thập niên 1980 và được hệ thống hóa trong công trình năm 1992
[8]. GP là một mở rộng trực tiếp của GA, tương ứng với dạng \textit{mã
hóa dựa trên cấu trúc} đã được giới thiệu trong
Mục~\ref{subsec:ma-hoa}. Khác biệt cốt lõi so với GA dùng mã hóa số
thực hay hoán vị là: thay vì tiến hóa một chuỗi giá trị có độ dài cố
định trong khuôn dạng nghiệm đã cho trước, GP tiến hóa trực tiếp
\textbf{cấu trúc} của một biểu thức hoặc chương trình. Nhờ đó, thuật
toán có khả năng tự khám phá đồng thời cả hình dạng lẫn tham số của
nghiệm -- đây chính là lý do GP được xem là bước chuyển từ ``tối ưu
tham số'' sang ``tự động tổng hợp chương trình''. Một tính chất đáng
chú ý của GP là nghiệm cuối cùng có dạng \textit{biểu thức tường minh},
có thể đọc, phân tích và giải thích được bởi con người, khác biệt căn
bản với các mô hình hộp đen như mạng nơ-ron.

\subsection*{Mã hóa dạng cây biểu thức}

Mỗi cá thể trong GP là một cây biểu thức. Mỗi nút trong của cây là một
toán tử được chọn từ một \textit{tập toán tử} cho trước, chẳng hạn các
phép toán số học $\{+,-,\times,\div\}$ hoặc hàm lượng giác
$\{\sin,\cos\}$; mỗi nút lá là một biến đầu vào hoặc hằng số được chọn
từ một \textit{tập biến và hằng số} cho trước. Khi duyệt cây theo đúng
cấu trúc phân cấp của nó, ta thu được một biểu thức toán học tường
minh, có thể tính giá trị trên bất kỳ bộ dữ liệu đầu vào nào. Điều này
cho phép toàn bộ quần thể các cây với hình dạng và nội dung khác nhau
được đánh giá, so sánh một cách thống nhất như trong GA thông thường.

Việc thiết kế hai tập hợp này cần thỏa mãn hai tính chất cơ bản:
\begin{itemize}
    \item \textbf{Tính khép kín} (closure): mọi toán tử phải chấp
    nhận bất kỳ giá trị nào do một toán tử khác hoặc một nút lá trả
    về, và luôn cho ra một kết quả hợp lệ. Ràng buộc này nảy sinh vì
    lai ghép và đột biến ghép các toán tử với nhau một cách ngẫu
    nhiên, không theo trật tự cố định nào. Chẳng hạn, nếu $\div$ xuất
    hiện tại một nút mà cây con ở vị trí mẫu số vừa tính ra $0$ --
    một tình huống hoàn toàn có thể xảy ra -- phép chia thông thường
    không xác định, khiến cả cây không tính được giá trị và cá thể
    đó trở nên vô nghĩa. Để tránh điều này, $\div$ thường được thay
    bằng \textit{phép chia bảo vệ}: trả về $0$ (hoặc một hằng số
    lớn) khi mẫu số bằng $0$, thay vì để phép tính thất bại. Các hàm
    có miền xác định hẹp hơn tập số thực như $\log$ hay $\sqrt{\cdot}$
    cũng cần phiên bản bảo vệ tương tự, chẳng hạn $\log(|x|)$ thay
    cho $\log(x)$, để luôn nhận được mọi giá trị đầu vào có thể xuất
    hiện trong cây.
    \item \textbf{Tính đủ} (sufficiency): hai tập hợp này phải chứa
    đủ các thành phần để biểu diễn được (ít nhất xấp xỉ) nghiệm của
    bài toán. Nếu dữ liệu thực sự mang quan hệ lượng giác mà tập toán
    tử chỉ gồm bốn phép tính số học, GP sẽ phải xấp xỉ hàm đích bằng
    đa thức với độ chính xác thấp.
\end{itemize}

Độ sâu tối đa của cây cũng là một tham số thiết kế quan trọng: giới hạn
quá chặt sẽ thu hẹp không gian khám phá, trong khi không giới hạn đủ
mức dễ dẫn đến các cây có kích thước rất lớn, tốn kém khi đánh giá và
khó diễn giải.

\textbf{Khởi tạo quần thể.} Khác với GA nơi quần thể ban đầu chỉ là các
chuỗi ngẫu nhiên, việc sinh cây ngẫu nhiên trong GP có ba chiến lược
kinh điển, tất cả đều giới hạn độ sâu tối đa $d_{\max}$:
\begin{itemize}
    \item \textbf{Full}: mọi nhánh của cây đều có
    độ sâu đúng bằng $d_{\max}$, tạo ra các cây cân đối, dày đặc;
    \item \textbf{Grow}: mỗi nhánh có thể kết thúc ở
    độ sâu bất kỳ không vượt quá $d_{\max}$, tạo ra các cây không đồng
    nhất về hình dạng;
    \item \textbf{Ramped half-and-half}:
    kết hợp hai phương pháp trên theo tỷ lệ $1:1$, đồng thời phân bố
    độ sâu tối đa của từng cây đều trên khoảng
    $[2, d_{\max}]$. Đây là chiến lược được khuyến nghị và sử dụng
    phổ biến nhất, vì nó tạo ra quần thể ban đầu đa dạng cả về kích
    thước lẫn hình dạng, giảm nguy cơ quần thể bị chi phối bởi một
    kiểu cấu trúc duy nhất ngay từ đầu.
\end{itemize}

\subsection*{Hàm độ thích nghi}

Ứng dụng tiêu biểu nhất của GP là bài toán \textbf{hồi quy ký hiệu}
(symbolic regression): tự động tìm một biểu thức toán học mô tả tốt
mối quan hệ giữa tập dữ liệu quan sát $(X, y)$ mà không giả định trước
dạng hàm. Bài toán này khác biệt căn bản so với hồi quy tuyến tính hay
hồi quy đa thức, nơi dạng hàm đã được ấn định sẵn và công việc còn lại
chỉ là tối ưu các hệ số. Trong GP, độ thích nghi của một cây thường
được xây dựng từ sai số dự đoán, tiêu biểu là sai số bình phương trung
bình giữa giá trị cây tính được và giá trị $y$ thực tế. Một số
công trình còn dùng RMSE, MAE, hoặc hệ số xác định $R^2$; lựa chọn
này ảnh hưởng trực tiếp đến độ nhạy của chọn lọc đối với các ngoại lai
trong dữ liệu.

Một vấn đề đặc trưng của GP là hiện tượng \textbf{bloat}: kích thước
cây có xu hướng tăng trưởng liên tục qua các thế hệ mà không mang lại
cải thiện tương xứng về độ chính xác. Nguyên nhân được lý giải ở mức
lý thuyết là các cây con trung tính (ví dụ đoạn $+0$, $\times 1$) tích
lũy dần vì chúng không làm giảm độ thích nghi nhưng lại được bảo vệ
khỏi đột biến có hại. Bloat làm tăng chi phí đánh giá, giảm khả năng
diễn giải nghiệm, và có thể dẫn đến quá khớp. Để đối phó, người ta
thường cộng thêm vào hàm thích nghi một \textbf{số hạng phạt} tỉ lệ với kích thước cây:
\begin{equation}
    \mathrm{fitness}(T) = \mathrm{MSE}(T) + \lambda \cdot |T|,
    \label{eq:parsimony}
\end{equation}
trong đó $|T|$ là số nút của cây $T$ và $\lambda > 0$ điều chỉnh mức
độ ưu tiên sự gọn nhẹ so với độ chính xác. Các biện pháp bổ sung gồm
giới hạn độ sâu/số nút tuyệt đối, loại bỏ cá thể vi phạm ngay sau khi
sinh, và các toán tử đột biến có chủ đích rút gọn cây.

\subsection*{Toán tử tiến hóa}

Các toán tử lai ghép và đột biến trong GP thao tác trực tiếp trên cấu
trúc cây thay vì trên chuỗi có độ dài cố định:
\begin{itemize}
    \item \textbf{Lai ghép trao đổi cây con}: chọn
    ngẫu nhiên một nút trên mỗi cây cha mẹ, sau đó hoán đổi hai cây con
    bắt rễ tại các nút đó cho nhau để tạo ra hai cây con mới. Các biến
    thể gồm lai ghép bất đối xứng (cây con thứ nhất lấy từ nút trong,
    cây con thứ hai từ nút lá) và lai ghép tự thân (cắt và dán trên
    cùng một cây), nhằm cân bằng giữa khám phá và khai thác;
    \item \textbf{Đột biến thay thế cây con}: chọn
    ngẫu nhiên một nút trên cây, rồi thay toàn bộ cây con bắt rễ tại nút
    đó bằng một cây con mới được sinh ngẫu nhiên;
    \item \textbf{Đột biến điểm} (point mutation): thay một nút hàm bằng
    một hàm khác cùng số ngôi, hoặc thay một nút lá bằng biến/hằng số
    khác, giữ nguyên cấu trúc cây;
    \item \textbf{Đột biến rút gọn} (hoist mutation): thay toàn bộ cây
    bằng một cây con ngẫu nhiên bên trong nó, qua đó thu nhỏ cây một
    cách có kiểm soát -- biến thể này thường được dùng kết hợp với số
    hạng phạt \eqref{eq:parsimony} để chống bloat.
\end{itemize}

Xác suất áp dụng của các toán tử cũng là tham số cần hiệu chỉnh: tỷ lệ
dành cho lai ghép, đột biến và sao chép trực tiếp (reproduction) cộng
lại bằng một, trong đó tỷ lệ lai ghép thường được đặt cao hơn hẳn tỷ
lệ đột biến; giá trị cụ thể tùy vào bài toán và được xác định bằng
thực nghiệm.

Các thành phần còn lại của chu trình tiến hóa, bao gồm khởi tạo quần
thể, chọn lọc, elitism và điều kiện dừng, về bản chất tương tự như GA
đã trình bày trong các mục trước; khác biệt duy nhất nằm ở đối tượng
thao tác là \textbf{cây} thay vì \textbf{vector} số hoặc hoán vị.
Chọn lọc giải đấu với kích thước giải đấu
$3$--$7$ được ưa chuộng trong GP vì đơn giản, ít nhạy với thang đo của
hàm thích nghi và dễ song song hóa khi kích thước quần thể lớn. Điều
kiện dừng thường là số thế hệ tối đa kết hợp với ngưỡng sai số chấp
nhận được; một số hệ thống còn theo dõi độ đa dạng của quần thể để kết
thúc sớm khi quần thể đã hội tụ.

\subsection*{Ví dụ minh họa}

Xét hai cây cha mẹ với tập toán tử $\{+,-,\times\}$. Cha mẹ thứ nhất
là cây $(x+1)^2$; do tập toán tử không có phép lũy thừa, cây này được
biểu diễn dưới dạng $\times[\,+[x,1],\ +[x,1]\,]$. Cha mẹ thứ hai là
cây $(x-1)x$, biểu diễn $\times[\,-[x,1],\ x\,]$, mang sẵn một
\textit{khối xây dựng} (building block) là nhân tử $(x-1)$.

Hình~\ref{fig:gp-crossover} minh họa phép lai ghép trao đổi cây con
giữa hai cha mẹ này. Nút cắt được chọn là nhánh $+[x,1]$ của cha mẹ
thứ nhất và nhánh $-[x,1]$ của cha mẹ thứ hai -- hai nhánh này được tô
màu cam trong hình. Hoán đổi hai nhánh cho nhau sinh ra hai cá thể
con: $(x-1)(x+1) = x^2 - 1$ và $(x+1)x = x^2 + x$, với phần vừa được
ghép sang từ cha mẹ kia tô màu xanh lá. Cả hai cá thể con đều tổ hợp
đúng nhân tử tuyến tính của cha mẹ này với phần còn lại của cha mẹ
kia, minh họa trực tiếp cách lai ghép cây con kết hợp các khối xây
dựng hữu ích từ nhiều cha mẹ -- cơ chế tương tự giả thuyết khối xây
dựng của GA nhưng vận hành trên cấu trúc phân cấp thay vì trên chuỗi
tuyến tính.

\begin{figure}[htbp]
\centering
\begin{tikzpicture}[
    op/.style={circle, draw, minimum size=7mm, inner sep=0pt, font=\small},
    leaf/.style={rectangle, draw, minimum size=6mm, inner sep=0pt, font=\small},
    cut/.style={fill=orange!25, draw=orange!70!black, thick},
    graft/.style={fill=green!20, draw=green!50!black, thick},
    level 1/.style={sibling distance=20mm, level distance=9mm},
    level 2/.style={sibling distance=9mm, level distance=9mm},
    edge from parent/.style={draw, -}
]

\begin{scope}
\node[op] {$\times$}
  child { node[op,cut] {$+$}
    child { node[leaf] {$x$} }
    child { node[leaf] {$1$} } }
  child { node[op] {$+$}
    child { node[leaf] {$x$} }
    child { node[leaf] {$1$} } };
\node[align=center, anchor=north] at (0,-2.3) {\small (a) Cha mẹ 1\\$(x+1)^2$};
\end{scope}

\begin{scope}[xshift=6cm]
\node[op] {$\times$}
  child { node[op,cut] {$-$}
    child { node[leaf] {$x$} }
    child { node[leaf] {$1$} } }
  child { node[leaf] {$x$} };
\node[align=center, anchor=north] at (0,-2.3) {\small (b) Cha mẹ 2\\$(x-1)x$};
\end{scope}

\node[align=center, font=\small] at (3,-3.9)
    {$\big\Downarrow$ Lai ghép trao đổi cây con};

\begin{scope}[yshift=-5.3cm]
\node[op] {$\times$}
  child { node[op,graft] {$-$}
    child { node[leaf] {$x$} }
    child { node[leaf] {$1$} } }
  child { node[op] {$+$}
    child { node[leaf] {$x$} }
    child { node[leaf] {$1$} } };
\node[align=center, anchor=north] at (0,-2.3) {\small (c) Con 1\\$(x-1)(x+1) = x^2-1$};
\end{scope}

\begin{scope}[xshift=6cm, yshift=-5.3cm]
\node[op] {$\times$}
  child { node[op,graft] {$+$}
    child { node[leaf] {$x$} }
    child { node[leaf] {$1$} } }
  child { node[leaf] {$x$} };
\node[align=center, anchor=north] at (0,-2.3) {\small (d) Con 2\\$(x+1)x = x^2+x$};
\end{scope}

\end{tikzpicture}
\caption{Lai ghép trao đổi cây con giữa hai cây cha mẹ (a)--(b), tạo
ra hai cây con (c)--(d). Nút cắt tô cam; phần vừa nhận từ cha mẹ kia
ở mỗi cây con tô xanh lá.}
\label{fig:gp-crossover}
\end{figure}

Ba toán tử đột biến cũng minh họa được trên một cây đơn giản, chẳng
hạn cây $x+3$, như trong Hình~\ref{fig:gp-mutation}. Đột biến thay thế
cây con chọn một nhánh -- ví dụ nhánh $3$ -- và thay bằng một cây con
mới sinh ngẫu nhiên, chẳng hạn $2x$, được cá thể $x+2x$. Đột biến điểm
giữ nguyên cấu trúc cây, chỉ đổi một nút hàm sang hàm khác cùng số
ngôi: đổi $+$ thành $-$, được cây $x-3$. Đột biến rút gọn thao tác
trên một cây lớn hơn, chẳng hạn $(x+3)\times x$: chọn cây con $x+3$
bên trong nó rồi thay thế toàn bộ cá thể bằng đúng cây con đó, được cá
thể mới chỉ còn $x+3$, nhỏ hơn hẳn cây gốc.

\begin{figure}[htbp]
\centering
\begin{tikzpicture}[
    op/.style={circle, draw, minimum size=7mm, inner sep=0pt, font=\small},
    leaf/.style={rectangle, draw, minimum size=6mm, inner sep=0pt, font=\small},
    cut/.style={fill=orange!25, draw=orange!70!black, thick},
    graft/.style={fill=green!20, draw=green!50!black, thick},
    level 1/.style={sibling distance=16mm, level distance=9mm},
    level 2/.style={sibling distance=9mm, level distance=9mm},
    edge from parent/.style={draw, -}
]

\begin{scope}
\node[op] {$+$}
  child { node[leaf] {$x$} }
  child { node[leaf,cut] {$3$} };
\node[align=center, anchor=north] at (0,-1.6) {\small (a) Trước\\$x+3$};
\end{scope}

\node[align=center, font=\footnotesize] at (2.6,-0.9)
    {$\Rightarrow$\\[1mm] thay $3$\\ bằng $2x$};

\begin{scope}[xshift=5.2cm]
\node[op] {$+$}
  child { node[leaf] {$x$} }
  child { node[op,graft] {$\times$}
    child { node[leaf] {$2$} }
    child { node[leaf] {$x$} } };
\node[align=center, anchor=north] at (0,-2.3) {\small (b) Sau\\$x+2x$};
\end{scope}

\node[align=center, font=\small\itshape] at (2.6,0.9) {Đột biến thay thế cây con};

\begin{scope}[yshift=-4.4cm]
\node[op,cut] {$+$}
  child { node[leaf] {$x$} }
  child { node[leaf] {$3$} };
\node[align=center, anchor=north] at (0,-1.6) {\small (c) Trước\\$x+3$};
\end{scope}

\node[align=center, font=\footnotesize] at (2.6,-5.3)
    {$\Rightarrow$\\[1mm] đổi $+$\\ thành $-$};

\begin{scope}[xshift=5.2cm, yshift=-4.4cm]
\node[op,graft] {$-$}
  child { node[leaf] {$x$} }
  child { node[leaf] {$3$} };
\node[align=center, anchor=north] at (0,-1.6) {\small (d) Sau\\$x-3$};
\end{scope}

\node[align=center, font=\small\itshape] at (2.6,-3.5) {Đột biến điểm};

\begin{scope}[yshift=-8.8cm]
\node[op] {$\times$}
  child { node[op,cut] {$+$}
    child { node[leaf] {$x$} }
    child { node[leaf] {$3$} } }
  child { node[leaf] {$x$} };
\node[align=center, anchor=north] at (0,-2.3) {\small (e) Trước\\$(x+3)\times x$};
\end{scope}

\node[align=center, font=\footnotesize] at (2.6,-9.7)
    {$\Rightarrow$\\[1mm] giữ lại\\ $x+3$};

\begin{scope}[xshift=5.2cm, yshift=-8.8cm]
\node[op,graft] {$+$}
  child { node[leaf] {$x$} }
  child { node[leaf] {$3$} };
\node[align=center, anchor=north] at (0,-1.6) {\small (f) Sau\\$x+3$};
\end{scope}

\node[align=center, font=\small\itshape] at (2.6,-7.9) {Đột biến rút gọn};

\end{tikzpicture}
\caption{Ba toán tử đột biến minh họa trên cây $x+3$ (đột biến rút
gọn minh họa trên cây lớn hơn $(x+3)\times x$). Nút hoặc nhánh bị thay
đổi tô cam trên cây trước đột biến (a, c, e); nút hoặc nhánh kết quả
tô xanh lá trên cây sau đột biến (b, d, f).}
\label{fig:gp-mutation}
\end{figure}

\newpage
\textbf{Minh họa bổ sung cho GA} (dành cho Mục~\ref{subsec:lai-ghep} và
Mục~\ref{subsec:dot-bien} của report chính -- xem trước ở đây trước khi
ghép vào đúng chỗ).

Hình~\ref{fig:ga-crossover} minh họa lai ghép đơn điểm trên mã hóa nhị
phân và lai ghép trộn trên mã hóa số thực -- hai cách mã hóa được nhắc
đến nhiều nhất trong report, lần lượt là cách mã hóa cổ điển của GA và
cách mã hóa dùng trong phần cài đặt thực nghiệm của tiểu luận. Với mã
hóa nhị phân, vị trí cắt được chọn ngẫu nhiên sau gen thứ $4$; đoạn từ
vị trí $5$ đến $8$ (tô cam) được hoán đổi giữa hai cha mẹ để tạo ra hai
cá thể con, với đoạn vừa nhận tô xanh lá. Với mã hóa số thực, hai cá
thể con được tính trực tiếp theo công thức lai ghép trộn
\eqref{eq:blend-crossover} với $\alpha = 0{,}75$.

\begin{figure}[H]
\centering
\begin{tikzpicture}[
    gene/.style={rectangle, draw, minimum width=7mm, minimum height=7mm, inner sep=0pt, font=\small},
    cut/.style={fill=orange!25, draw=orange!70!black, thick},
    graft/.style={fill=green!20, draw=green!50!black, thick}
]

\node[font=\small\itshape] at (2.8,1.0) {Mã hóa nhị phân -- lai ghép đơn điểm};

\node[gene] at (0.0,0) {1};
\node[gene] at (0.8,0) {0};
\node[gene] at (1.6,0) {1};
\node[gene] at (2.4,0) {1};
\node[gene,cut] at (3.2,0) {0};
\node[gene,cut] at (4.0,0) {0};
\node[gene,cut] at (4.8,0) {1};
\node[gene,cut] at (5.6,0) {0};
\node[anchor=east, font=\small] at (-0.5,0) {Cha mẹ 1};

\node[gene] at (0.0,-0.9) {0};
\node[gene] at (0.8,-0.9) {1};
\node[gene] at (1.6,-0.9) {1};
\node[gene] at (2.4,-0.9) {0};
\node[gene,cut] at (3.2,-0.9) {1};
\node[gene,cut] at (4.0,-0.9) {0};
\node[gene,cut] at (4.8,-0.9) {1};
\node[gene,cut] at (5.6,-0.9) {1};
\node[anchor=east, font=\small] at (-0.5,-0.9) {Cha mẹ 2};

\draw[dashed] (2.8,0.5) -- (2.8,-1.4);

\node[align=center, font=\small] at (2.8,-1.9) {$\Downarrow$ hoán đổi đoạn sau vị trí cắt};

\node[gene] at (0.0,-2.8) {1};
\node[gene] at (0.8,-2.8) {0};
\node[gene] at (1.6,-2.8) {1};
\node[gene] at (2.4,-2.8) {1};
\node[gene,graft] at (3.2,-2.8) {1};
\node[gene,graft] at (4.0,-2.8) {0};
\node[gene,graft] at (4.8,-2.8) {1};
\node[gene,graft] at (5.6,-2.8) {1};
\node[anchor=east, font=\small] at (-0.5,-2.8) {Con 1};

\node[gene] at (0.0,-3.7) {0};
\node[gene] at (0.8,-3.7) {1};
\node[gene] at (1.6,-3.7) {1};
\node[gene] at (2.4,-3.7) {0};
\node[gene,graft] at (3.2,-3.7) {0};
\node[gene,graft] at (4.0,-3.7) {0};
\node[gene,graft] at (4.8,-3.7) {1};
\node[gene,graft] at (5.6,-3.7) {0};
\node[anchor=east, font=\small] at (-0.5,-3.7) {Con 2};

\node[font=\small\itshape] at (2.8,-5.0) {Mã hóa số thực -- lai ghép trộn ($\alpha = 0{,}75$)};

\node[gene, minimum width=13mm] at (0.5,-5.9) {$3{,}0$};
\node[gene, minimum width=13mm] at (2.15,-5.9) {$-1{,}0$};
\node[anchor=east, font=\small] at (-0.9,-5.9) {Cha mẹ 1};

\node[gene, minimum width=13mm] at (0.5,-6.8) {$1{,}0$};
\node[gene, minimum width=13mm] at (2.15,-6.8) {$2{,}0$};
\node[anchor=east, font=\small] at (-0.9,-6.8) {Cha mẹ 2};

\node[align=center, font=\small] at (1.3,-7.3) {$\Downarrow$ \eqref{eq:blend-crossover}};

\node[gene, minimum width=13mm] at (0.5,-8.2) {$2{,}5$};
\node[gene, minimum width=13mm] at (2.15,-8.2) {$-0{,}25$};
\node[anchor=east, font=\small] at (-0.9,-8.2) {Con 1};

\node[gene, minimum width=13mm] at (0.5,-9.1) {$1{,}5$};
\node[gene, minimum width=13mm] at (2.15,-9.1) {$1{,}25$};
\node[anchor=east, font=\small] at (-0.9,-9.1) {Con 2};

\end{tikzpicture}
\caption{Ví dụ lai ghép trong GA trên hai cách mã hóa. Trên: mã hóa
nhị phân với lai ghép đơn điểm -- đoạn sắp hoán đổi tô cam trên cha
mẹ, đoạn vừa nhận được ở cá thể con tô xanh lá. Dưới: mã hóa số thực
với lai ghép trộn theo \eqref{eq:blend-crossover}, $\alpha = 0{,}75$.}
\label{fig:ga-crossover}
\end{figure}

Tương tự, Hình~\ref{fig:ga-mutation} minh họa đột biến đảo bit trên mã
hóa nhị phân và đột biến Gauss trên mã hóa số thực. Với mã hóa nhị
phân, gen tại vị trí $3$ (tô cam) được đảo giá trị từ $1$ thành $0$.
Với mã hóa số thực, giá trị tại vị trí $2$ (tô cam) được cộng thêm một
nhiễu Gauss theo \eqref{eq:dot-bien-gauss}, ở đây lấy mẫu được
$\mathcal{N}(0,\sigma^2) = -0{,}4$, cho ra giá trị mới $-1{,}4$ (tô
xanh lá).

\begin{figure}[H]
\centering
\begin{tikzpicture}[
    gene/.style={rectangle, draw, minimum width=7mm, minimum height=7mm, inner sep=0pt, font=\small},
    cut/.style={fill=orange!25, draw=orange!70!black, thick},
    graft/.style={fill=green!20, draw=green!50!black, thick}
]

\node[font=\small\itshape] at (2.8,1.0) {Mã hóa nhị phân -- đột biến đảo bit};

\node[gene] at (0.0,0) {1};
\node[gene] at (0.8,0) {0};
\node[gene,cut] at (1.6,0) {1};
\node[gene] at (2.4,0) {1};
\node[gene] at (3.2,0) {0};
\node[gene] at (4.0,0) {0};
\node[gene] at (4.8,0) {1};
\node[gene] at (5.6,0) {0};
\node[anchor=east, font=\small] at (-0.5,0) {Trước};

\node[align=center, font=\small] at (2.8,-0.7) {$\Downarrow$ đảo bit tại vị trí 3};

\node[gene] at (0.0,-1.6) {1};
\node[gene] at (0.8,-1.6) {0};
\node[gene,graft] at (1.6,-1.6) {0};
\node[gene] at (2.4,-1.6) {1};
\node[gene] at (3.2,-1.6) {0};
\node[gene] at (4.0,-1.6) {0};
\node[gene] at (4.8,-1.6) {1};
\node[gene] at (5.6,-1.6) {0};
\node[anchor=east, font=\small] at (-0.5,-1.6) {Sau};

\node[font=\small\itshape] at (2.8,-2.8) {Mã hóa số thực -- đột biến Gauss};

\node[gene, minimum width=13mm] at (0.5,-3.7) {$3{,}0$};
\node[gene, cut, minimum width=13mm] at (2.15,-3.7) {$-1{,}0$};
\node[anchor=east, font=\small] at (-0.9,-3.7) {Trước};

\node[align=center, font=\small] at (1.3,-4.4) {$\Downarrow$ cộng nhiễu $\mathcal{N}(0,\sigma^2)$ vào vị trí 2};

\node[gene, minimum width=13mm] at (0.5,-5.3) {$3{,}0$};
\node[gene, graft, minimum width=13mm] at (2.15,-5.3) {$-1{,}4$};
\node[anchor=east, font=\small] at (-0.9,-5.3) {Sau};

\end{tikzpicture}
\caption{Ví dụ đột biến trong GA trên hai cách mã hóa. Trên: mã hóa
nhị phân với đột biến đảo bit tại vị trí thứ $3$. Dưới: mã hóa số
thực với đột biến Gauss, cộng nhiễu ngẫu nhiên vào vị trí thứ $2$
theo \eqref{eq:dot-bien-gauss}. Vị trí bị thay đổi tô cam trước đột
biến; giá trị kết quả tô xanh lá sau đột biến.}
\label{fig:ga-mutation}
\end{figure}
